\section{Reproducibility Specification}

This section states the implementation contract independently of any particular
storage layout. A researcher can implement the method in any programming
environment if the implementation preserves the inputs, artifacts, normalization
rules, and validation checks described below.

\subsection{Run Configuration}

Each run is defined by a frozen configuration. The configuration must include
the legal source text or source identifier, a run identifier, the response-model
roster, sample count per model, stage-specific temperature policy,
perturbation policy, clustering method, judge model, quality gates, and output
location. The temperature policy must distinguish deterministic benchmark
construction from stochastic response generation: Frank and Karthic may be
fixed at temperature 0 for reproducible source and rubric artifacts, while
response models should use nonzero sampling temperatures and repeated samples
when the goal is to observe natural variation in legal reasoning. Judge calls
should record their temperature and, for reliability claims, use repeats or a
panel so score variability is measured. The configuration should also record
whether the run is offline, live, or a fixture stress test. All stochastic
choices, including model temperature offsets and fixture sampling seeds, must be
explicit.

The reference implementation materializes this contract with a protocol-freeze
manifest before a run is treated as frozen. The freeze manifest records the
configuration hash, source hash, ordered instruction-stream hashes for Frank,
Karthic, and Dasha, response-model roster, clustering settings, judge model or
judge-panel roster, perturbation policy, quality gates, normalization versions,
and a top-level protocol hash. It intentionally omits API keys and generated
model outputs. This separates the research protocol from mutable local
environment state and makes later result bundles auditable against the exact
method that was intended to run.

The repository also includes a static manuscript linter for environments without
a TeX engine. The linter does not replace \texttt{latexmk} or
\texttt{pdflatex}, but it verifies that all \texttt{\textbackslash input}
targets, figure files, bibliography files, and citation keys referenced by the
paper resolve. A compiled PDF remains the stronger artifact when a full TeX
installation is available.

Finally, the implementation provides a no-call preflight check for live
configs. Preflight builds the protocol-freeze metadata in memory and checks that
the run is live-capable, uses natural response prompting, has a sufficiently
diverse response roster, uses LLM agents for Frank, Karthic, and Dasha, uses
reasoning-signature clustering, requires multiple observed clusters, configures
LLM judging with repeats or a panel, and can find the source case. It also
reports local credential and perturbation-policy warnings before any model call
is made. This makes budgeted live runs auditable before execution rather than
only after artifacts have been generated.

\begin{table}[h]
\centering
\caption{Minimum run-configuration fields required to reproduce a pipeline run.}
\label{tab:run-config-contract}
\begin{tabular}{p{0.24\linewidth}p{0.66\linewidth}}
\toprule
Field group & Required contents \\
\midrule
Source & Source-case text or stable source identifier; jurisdiction if the benchmark depends on jurisdiction-specific law; source hash. \\
Agents & Model identifier, provider adapter, temperature, instruction-stream hash, and output schema for Frank, Karthic, Dasha extraction, and Judge. \\
Responses & Response-model identifiers, sample count per model, temperature policy, response prompt style, and question track ids. \\
Perturbations & Whether perturbations are enabled; maximum number of variations; requirement for invariant and material tracks; expected behavior for each track. \\
Clustering & Signature schema version; normalization version; minimum observed clusters; mixed-cluster threshold; representative-selection rule. \\
Judge and Zak & Rubric scale, judge aggregation rule, disagreement threshold, escalation margin, and packet-generation rule. \\
Outputs & Versioned JSON artifacts, prompt hashes, response provenance, validation summary, generated tables, and figure inputs. \\
\bottomrule
\end{tabular}
\end{table}

\subsection{Artifact Schemas}

The pipeline is reproducible because stages exchange versioned artifacts rather
than conversational state. Frank emits the benchmark packet. Karthic emits
rubric rows. Response generation emits raw model answers with question-track
metadata. Dasha emits clusters with member ids and legal signals. Judge emits
row-level centroid scores and projected member scores. Zak emits either an empty
no-escalation state or an escalation packet.

\begin{table}[h]
\centering
\caption{Required artifact interfaces for an independent implementation.}
\label{tab:artifact-contract}
\begin{tabular}{p{0.18\linewidth}p{0.72\linewidth}}
\toprule
Artifact & Required fields \\
\midrule
Frank packet & \texttt{id}, \texttt{source\_hash}, \texttt{doctrine\_family}, \texttt{detected\_doctrine\_gates}, \texttt{source\_extraction}, \texttt{neutral\_question}, \texttt{gold\_answer}, \texttt{variations}, \texttt{controller\_card}, \texttt{prompt\_hashes}. \\
Karthic rubric & \texttt{id}, \texttt{source\_hash}, \texttt{frank\_packet\_id}, \texttt{rows}, \texttt{weights}, \texttt{prompt\_hashes}. Each row requires an id, category, weight, criterion, source-support field, and scoring anchors or equivalent rubric guidance. \\
Responses & \texttt{id}, \texttt{model}, \texttt{sample\_index}, \texttt{text}, \texttt{question\_id}, \texttt{track\_id}, \texttt{variant\_id}, \texttt{perturbation\_type}, \texttt{expected\_behavior}, timestamp or generation provenance. \\
Dasha clusters & \texttt{id}, \texttt{method}, \texttt{legal\_signal}, \texttt{representative\_response\_id}, \texttt{member\_response\_ids}, \texttt{size}, \texttt{centroid\_quality}, and, in perturbation runs, question-track metadata. \\
Judge scores & One row-level score and rationale for every rubric row and every cluster representative; weighted centroid scores; projected member scores; model rankings. \\
Zak packets & Escalation status, triggering threshold, disputed cluster ids, disputed rubric row ids, rationale, and handoff text when escalation is required. \\
\bottomrule
\end{tabular}
\end{table}

\subsection{Instruction-Stream Construction}

The agent context is an ordered prompt bundle, not hidden memory. A reproducible
implementation should store each prompt component with a stable component id,
role, version, and text. At runtime, the implementation concatenates the
components in fixed order, prefixes each component with its id, truncates only at
a declared token or character budget, and records a cryptographic hash over the
rendered stream. The same rendered stream and upstream artifact must be supplied
to the same model identifier to reproduce the stage input.

The shared first component is a cross-agent legal reasoning protocol. It states
that source authority controls legal substance, that calibration examples do not
override the source, that schema fields are doctrine-general, that uncertainty is
preserved rather than silently filled, and that perturbations must identify the
expected legal effect of the edited fact. Role-specific components then define
the agent contract. Frank receives source-intake, packet-building, question
writing, gold-answer, and variation instructions. Karthic receives packet-audit,
row-card, weighting, duplicate-control, and scoring-overlay instructions. Dasha
receives signature-extraction, track-audit, clustering, centroid-metadata, and
escalation instructions. Judge receives rubric-application and row-score
rationale instructions.

\subsection{Agent Calls}

Frank receives the source text and its instruction stream, then returns one
strict JSON packet. Karthic receives the locked Frank packet and its instruction
stream, then returns one strict JSON rubric. Response models receive only the
question text for the active track under the natural response protocol; they do
not receive the gold answer, rubric, hidden reasoning labels, or required memo
headings. Dasha receives the locked Frank packet, its instruction stream, and one
raw response at a time, then returns a strict JSON reasoning signature. Judge
receives one Dasha representative, the cluster legal signal, and all Karthic row
criteria, then returns one score and rationale per row.

\subsection{Dasha Clustering Algorithm}

The reported configuration clusters responses by Dasha-extracted
legal-reasoning signatures rather than by surface text. For each response,
Dasha emits both natural-language fields and canonical identifiers:
\texttt{doctrine\_id}, \texttt{rule\_trigger\_id}, \texttt{outcome\_id},
\texttt{exception\_or\_defense\_id}, and \texttt{primary\_reasoning\_id}. The
identifier vocabulary is not pre-enumerated in code. Dasha derives it from the
Frank packet, source-derived gates, instruction context, and response text.
The implementation sanitizes the signature, removes unexpected script
fragments, collapses whitespace, and groups exact canonical identifiers.

The cluster key is the tuple
\texttt{(doctrine\_id, rule\_trigger\_id, outcome\_id,
exception\_or\_defense\_id, primary\_reasoning\_id,
secondary\_cluster\_profile)}. The first five values are Dasha's canonical
legal abstraction. The final value is derived from Dasha's secondary-path list
and includes only accepted or uncertain non-primary paths. Rejected and merely
mentioned side paths remain in the audit profile but do not automatically split
clusters. Legacy keyword buckets exist only for archived fixtures without
Dasha canonical identifiers.

Responses are grouped by exact cluster-key match within each question track. A
perturbation run never clusters different tracks together. For each cluster, the
representative response is the member with the highest average similarity to the
other members, combining token-level Jaccard similarity with feature agreement
over gate, outcome, exception, and reasoning fields. Text similarity is retained
as a centroid-quality diagnostic, not as the clustering key. A cluster is flagged
for review when sampled members visibly use different legal theories, when the
normalized signal conflicts with the response text, or when centroid/member
feature agreement falls below the configured threshold.

\subsection{Score Projection and Ranking}

The judge scores only cluster representatives. Each representative receives a
0--4 score and rationale for every Karthic rubric row. The weighted centroid
score is then copied to every response assigned to that centroid. Model rankings
are computed from the projected member scores, typically by averaging all
projected scores for a model across samples and question tracks. This projection
rule is valid only if Dasha cluster membership is legally coherent; therefore
cluster audit and perturbation validation are prerequisites for interpreting
rankings.
